\section{Introduction}
\label{sec:introduction}

Electrofused grains, notably silicon carbide and fused aluminum oxide,
are synthetic raw materials employed in the manufacture of abrasive,
refractory, and ceramic products. Production starts in
electric-resistance (Acheson-type) or electric-arc furnaces, whose
output blocks are subsequently crushed, milled, and classified by sets
of screens into a large number of granulometric ranges (grits) with
strict physical--chemical specifications. A distinctive feature of the
grain-finishing stage is \emph{co-production in fixed proportions}:
each configuration of the crushing and classification equipment --
hereafter a \emph{production process} -- simultaneously yields a
specific mix of final products in technologically determined
proportions. Planners therefore do not decide production quantities
product by product; they decide \emph{which process to operate in each
period}, and the product mix follows. Combined with lengthy setups
(equipment reconfiguration and cleaning may consume an entire shift),
a portfolio of dozens of products, and tight order-book due dates,
this makes production scheduling a genuinely hard combinatorial task.

The problem was formalized by \citet{luche2009combining} as a
\emph{process selection and lot-sizing problem} (PSP): a single-stage,
big-bucket model in which at most one process may be assigned to each
period and the objective is to minimize the total shortage of the
order book (with a small penalty on excess production). At the time,
the mixed-integer programming (MIP) model could not be solved to
optimality within acceptable times for instances of practical size by
the then state-of-the-art solver, which motivated constructive
heuristics, local search, and later a GRASP procedure
\todo{cite thesis}. Two developments justify revisiting the problem
fifteen years later. First, MIP solvers have improved by orders of
magnitude in this period through advances in presolve, cut separation,
primal heuristics, and parallel tree search
\citep{koch2022progress,achterberg2013analyzing,achterberg2020presolve};
whether such progress has effectively closed a practical problem that
was intractable in 2011 is an empirical question of direct managerial
relevance. Second, MIP-based construction and improvement heuristics
-- \emph{relax-and-fix} (\RF) and \emph{fix-and-optimize} (\FO), now
standard components of the matheuristics toolbox
\citep{sahling2009solving,helber2010fix,toledo2015relax} -- have
matured precisely on lot-sizing structures similar to the PSP, but
have never been applied to it.

This paper reports a comprehensive computational study of the PSP that
confronts three algorithmic paradigms under a strict equal-resources
protocol: (i) the monolithic MIP model solved by a modern commercial
solver; (ii) temporal-decomposition matheuristics (\RFFO{} and an
\RF-warm-started monolithic variant, \RFMIP); and (iii) a
state-of-the-art hybrid metaheuristic combining reactive GRASP,
iterated local search, variable neighborhood descent, tabu components,
and cross-run path relinking (\ILS). All methods run on the same
machine, with the same wall-clock budgets and explicit thread
accounting, and the central comparisons were specified as
pre-registered hypotheses before the experiments were executed.
Rather than asking only \emph{whether} a heuristic beats the solver,
we map \emph{when} each method is the tool of choice as a function of
instance scale (planning horizons from $T\!=\!19$ to $T\!\approx\!190$
periods) and available computing time (from operational re-planning
budgets of 5--10 minutes to tactical budgets of 1--3 hours).
\sprintii{one-sentence preview of the frontier result once Sprint~2
hypotheses H1--H2 are adjudicated.}

The contributions are fourfold. First, we quantify fifteen years of
MIP solver progress on a documented real-world problem: instances that
exhibited large optimality gaps in 2011 are now solved in seconds
\sprintii{exact figures}, while a hard subclass remains open at the
largest scales -- a data point for the solver-progress literature
\citep{koch2022progress} grounded in an industrial application.
Second, we develop and specify, at a reproducibility level rarely
reported (bound-update mechanics, per-window time-limit policies, and
budget guards included), \RF{} and \FO{} matheuristics tailored to the
co-production structure of the PSP, together with an \RFMIP{} hybrid
that dominates the cold-started solver by construction. Third, we
provide what is, to the best of our knowledge, the first publicly
available benchmark for lot sizing with co-production in fixed
proportions: 60 instances spanning ten problem scales, with best-known
solutions, dual bounds, and full per-run trajectories, released in
open formats. Fourth, we derive managerial insights for the
electrofused grains mill, including the practical consequence that
order-book scheduling can now be re-optimized within the minutes-scale
budgets of daily re-planning \sprintii{support with H1 outcome}.

The remainder of the paper is organized as follows.
Section~\ref{sec:related} reviews the related literature and positions
the PSP within the lot-sizing and scheduling taxonomy.
Section~\ref{sec:problem} describes the industrial setting and the
mathematical model. Section~\ref{sec:matheuristics} presents the
\RF{}, \FO{}, and \RFMIP{} methods, and
Section~\ref{sec:metaheuristic} summarizes the hybrid \ILS{} baseline.
Section~\ref{sec:experiments} reports the experimental protocol and
results, and Section~\ref{sec:conclusions} concludes.
